# Tämä oli kappaleessa 4.3 Putket
\begin{maaritelma}
    Olkoon $f:\C^2\to\C$ polynomikuvaus. Sen kuitua $f^{-1}(c),c\in\C$ kutsutaan \emph{säännölliseksi}, mikäli jollekin $c$:n ympärisölle $D$ kuvaus $f|_{f^{-1}(D)}:f^{-1}(D)\to D$ on paikallisesti triviaali $C^\infty$-kuidutus.

    Kuitu $f^{-1}(c)$ on \emph{säännöllinen äärettömyydessä}, jos jollekin $c$:n ympäristölle $E$ ja kompaktille osajoukolle $K\subset\C^2$ kuvaus $f|_{f^{-1}(D)-K}$ on paikallisesti triviaali kuidutus. 

    Jos kaikki $f$:n kuidut ovat säännöllisiä äärettömyydessä niin $f$ on \emph{hyvä}.
\end{maaritelma}

# Epäonnistunut homologiaosio
\subsection{Homologia (aksiomaattinen yritys)}
Tämä on vähän epätoivoinen yritys sisällyttää homologia työhön. Poistetaan tämä ja laitetaan ylempään homologiaosioon pari määritelmää ja lausetta.

Tämän aliosion tarkoitus on kiinnittää työssä käytettävä homologian sanasto. Työn tiivis esitys ei tee oikeutta homologialle ja homologiaa tuntematonta lukijaa suositellaankin sivuuttamaan osio, tai tutustumaan homologiaan Hatcherin kirjan avulla \citep{Hatcher2002}. Aliosio on pikemminkin suunnattu homologiaa tuntevalle lukijalle, joka haluaa varmistaa termistön, tai palata työn myöhemmistä osioista tarkistamaan määritelmiä. Sanastossa noudatetaan kuitenkin yleisiä tapoja, joten myös homologiaa tunteva lukija voi sivuuttaa aliosion.
\begin{maaritelma}
    \emph{CW-kompleksi} on topologinen avaruus $X$, joka muodostetaan algoritmisesti:
    \begin{enumerate}
        \item Valitaan diskreetti joukko $X^0$, jonka alkioita kutsutaan $X$:n \emph{$0$-soluiksi}.
        \item Muodostetaan induktiivisesti \emph{$n$-luuranko $X^n$} $X^{n-1}$:stä liittämällä $n$-soluja $e_\alpha^n$ kuvauksilla $\varphi_\alpha:\pallo^{n-1}\to X^{n-1}$. $X^n$ on tekijä-avaruus
        \begin{equation}
            X^n=\frac{X^{n-1}\sqcup_\alpha D_\alpha^n}{\sim},
        \end{equation}
        missä $\sim$ on relaatio $x\sim\varphi_\alpha(x)$, kun $x\in\partial D_\alpha^n$. Joukon $D_\alpha^n-\partial D_\alpha^n$ kuva tekijäkuvauksen suhteen ja solu $e_\alpha^n$ ja ovat homeomorfisia.
        \item Jatkamalla prosessia saadaan $X=\cup_n X^n$ ja se voidaan varustaa heikolla topologialla: $A\subset X$ on avoin (suljettu) jos ja vain jos $A\cap X^n$ on avoin (suljettu) $X^n$:ssä kaikille $n\in\N$. 
    \end{enumerate}
\end{maaritelma}
\begin{aksiooma}[Homologia-aksioomat]
    \begin{enumerate}
        \item Jos $f\cong g:X\to Y$ niin $f_\bullet=g_\bullet:\tilde h_n(X)\to\tilde h_n(Y)$.
        \item On olemassa reunahomomorfismit $\partial:\tilde h_n(X/A)\to\tilde h_{n-1}(A)$ määritelty kaikille CW-pareille $(X,A)$, siten että jono
        \begin{equation}
            \cdots\xrightarrow{\partial}\tilde h_n(A)\xrightarrow{i_\bullet}\xrightarrow{q_\bullet}\tilde h_n(X/A)\xrightarrow{\partial}\tilde h_n(A)\xrightarrow{i_\bullet}\cdots
        \end{equation}
        on tarkka. Jonossa $i$ on inkluusio ja $q$ tekijäkuvaus.
    \end{enumerate}
\end{aksiooma}
\begin{maaritelma}
    Pelkistetty homologiateoria antaa epätyhjille CW-komplekseille $X$ jonon
    \begin{itemize}
        \item abelisia ryhmiä $\tilde h_n(X)$ ja
        \item homomorfismeja $f_\bullet:\tilde h_n(X)\to\tilde h_n(Y)$ jokaista CW-kompleksien välistä kuvausta $f:X\to Y$ kohden, siten että
        \begin{itemize}
            \item $(fg)_\bullet=f_\bullet g_\bullet$
            \item $\mathbbm{1}_\bullet=\mathbbm{1}$,
        \end{itemize}
    \end{itemize}
    siten, että homologia-aksioomat pitävät.
\end{maaritelma}

# Linkkiluku
\begin{maaritelma}[Määritelmä 1: pidän tätä varmuuden vuoksi, mutta ajattelin poistaa]
    Olkoot $\gamma_1,\gamma_2:\S^1\to\R^3$ eli leikkaamattomia sileitä käyriä. Tällöin käyrien \emph{linkkiluku} on
    \begin{equation*}
        \ell(\gamma_1,\gamma_2):=\frac{1}{4\pi}\int_{\S^1\times\S^1}\frac{\det(\dot{\gamma_1}(s),\dot{\gamma_2}(t),\gamma_1(s)-\gamma_2(t))}{|\gamma_1(s)-\gamma_2(t)|^3}.
    \end{equation*}
    Riittäisikö tähän vähän epämuodollisempi määrittely, joka perustuu kuviin? Toi integraali on aika tuhti.
\end{maaritelma}